%% =====================================================================
%% 【示例 / SYNTHETIC EXAMPLE】example-mechanism.tex
%% This is a SYNTHETIC toy pathway for demonstrating the biopathway.tex
%% style library ONLY. All entities/edges below are invented; it is NOT
%% derived from any real case material and must never be shipped as a
%% case deliverable. 本文件为合成示例，非真实案例结果。
%%
%% Compile:  pdflatex example-mechanism.tex   (needs tikz + pgfplots)
%% Entities: LIGAND-X (small molecule), RcptA (receptor), KinB (protein),
%%           TFc (gene/TF), TargetD (gene). 3 edge types shown:
%%           activate / inhibit / phosphorylate (+ transport).
%% =====================================================================
\documentclass[tikz,border=6pt]{standalone}
%% =====================================================================
%% biopathway.tex -- TikZ/pgfplots style library for biological mechanism
%% and pathway diagrams (生物机制图与通路图 图元库).
%%
%% Usage in a figure file:
%%   \documentclass[tikz,border=6pt]{standalone}
%%   %% =====================================================================
%% biopathway.tex -- TikZ/pgfplots style library for biological mechanism
%% and pathway diagrams (生物机制图与通路图 图元库).
%%
%% Usage in a figure file:
%%   \documentclass[tikz,border=6pt]{standalone}
%%   %% =====================================================================
%% biopathway.tex -- TikZ/pgfplots style library for biological mechanism
%% and pathway diagrams (生物机制图与通路图 图元库).
%%
%% Usage in a figure file:
%%   \documentclass[tikz,border=6pt]{standalone}
%%   \input{biopathway.tex}        % this file (same directory or TEXINPUTS)
%%   \begin{document} ... \end{document}
%%
%% Dependencies (deliberately minimal): tikz, pgfplots, xcolor only.
%% TikZ libraries used: arrows.meta, positioning, calc, backgrounds,
%% fit, decorations.pathmorphing -- all ship with every full TeX Live.
%% Compiles with pdflatex (also xelatex). No exotic packages.
%% =====================================================================
\RequirePackage{xcolor}
\RequirePackage{tikz}
\usetikzlibrary{arrows.meta,positioning,calc,backgrounds,fit,decorations.pathmorphing}
\RequirePackage{pgfplots}
\pgfplotsset{compat=1.17} % safe baseline; 1.16 also fine if older TeX Live

%% ---------- palette (colorblind-aware, Okabe-Ito inspired) ----------
\definecolor{bpblue}{HTML}{0072B2}    % protein / receptor fill
\definecolor{bpblueL}{HTML}{BFDFF2}   % protein light fill
\definecolor{bporange}{HTML}{E69F00}  % gene / nucleic fill
\definecolor{bpgreen}{HTML}{009E73}   % small molecule
\definecolor{bppurple}{HTML}{CC79A7}  % complex
\definecolor{bpred}{HTML}{D55E00}     % inhibition / danger
\definecolor{bpgray}{HTML}{999999}    % uncertain
\definecolor{membrane}{HTML}{F2E3C9}  % lipid bilayer band
\definecolor{bpbg}{HTML}{F7F9FB}      % compartment background tint

%% ---------- compartment bands ----------
%% Horizontal bands, top->bottom: extracellular, membrane, cytosol, nucleus.
%% Draw them on the background layer FIRST, then place nodes on top.
%% Usage: \fill[extracellularband] (-8,2.5) rectangle (8,5);
\tikzset{
  extracellularband/.style={fill=white},
  membraneband/.style={fill=membrane},
  cytosolband/.style={fill=bpbg},
  nucleusband/.style={fill=bpblueL!35},
  bandlabel/.style={font=\small\itshape, text=black!55, anchor=north west},
}

%% ---------- membrane bilayer pic ----------
%% A double line of lipid head circles + tails spanning a horizontal band.
%% Usage: \pic at (0,0) {bilayer={8}};  % half-width 8cm
\tikzset{
  pics/bilayer/.style={
    code={
      \pgfmathsetmacro{\hw}{#1}
      \begin{scope}
        \foreach \i in {0,...,40} {
          \pgfmathsetmacro{\px}{-\hw + \i*(2*\hw/40)}
          \draw[fill=bpred!70, draw=bpred!70] (\px,0.28) circle (0.09);
          \draw[fill=bpred!70, draw=bpred!70] (\px,-0.28) circle (0.09);
          \draw[bporange!80, line width=0.7pt] (\px,0.19) -- (\px,0.03);
          \draw[bporange!80, line width=0.7pt] (\px,-0.19) -- (\px,-0.03);
        }
      \end{scope}
    }
  },
}

%% ---------- node styles ----------
\tikzset{
  %% protein: rounded rectangle, upright (正体) label
  protein/.style={rectangle, rounded corners=3pt, draw=bpblue,
    fill=bpblueL, thick, minimum width=1.7cm, minimum height=0.8cm,
    align=center, font=\small},
  %% gene: italic (斜体) label, orange accent. Use \gene{IL6} macro below
  %% to guarantee italic math-style name.
  gene/.style={rectangle, draw=bporange, fill=bporange!18, thick,
    minimum width=1.4cm, minimum height=0.7cm, align=center,
    font=\small\itshape},
  %% complex: two stacked rounded boxes look (purple)
  complex/.style={rectangle, rounded corners=6pt, draw=bppurple,
    fill=bppurple!20, thick, double, double distance=1.2pt,
    minimum width=2.0cm, minimum height=0.9cm, align=center, font=\small},
  %% small molecule / metabolite: ellipse, green
  smallmolecule/.style={ellipse, draw=bpgreen, fill=bpgreen!15, thick,
    minimum width=1.3cm, minimum height=0.7cm, align=center, font=\small},
  %% transmembrane receptor: tall rounded rect meant to SPAN the membrane
  %% band (place its center on the membrane midline).
  receptor/.style={rectangle, rounded corners=8pt, draw=bpblue,
    fill=bpblueL, very thick, minimum width=1.1cm, minimum height=2.0cm,
    align=center, font=\small},
  %% nucleus body (large double circle) for a nuclear compartment
  nucleusbody/.style={circle, draw=bpblue!60, fill=nucleusband, thick,
    double, double distance=2pt, minimum size=5.2cm},
  %% small badges: phosphate / ubiquitin
  pbadge/.style={circle, draw=bpred, fill=white, inner sep=0.6pt,
    minimum size=0.42cm, font=\scriptsize\bfseries},
  ubbadge/.style={circle, draw=bppurple, fill=white, inner sep=0.6pt,
    minimum size=0.55cm, font=\scriptsize\bfseries},
  %% generic annotation text
  ann/.style={font=\scriptsize, text=black!70, align=center},
}

%% Italic gene name helper: \gene{STAT3} -> italic STAT3
\newcommand{\gene}[1]{\textit{#1}}
%% Phosphorylated protein: \phos{STAT3} -> p-STAT3 (upright p- prefix)
\newcommand{\phos}[1]{p-#1}

%% ---------- edge styles ----------
%% arrows.meta tips; every edge carries its mechanism semantics in style.
\tikzset{
  %% activation 促进/激活: open triangle arrow
  activate/.style={-{Triangle[open,length=7pt]}, thick, draw=bpblue,
    shorten >=1pt, shorten <=1pt},
  %% inhibition 抑制: perpendicular T-bar
  inhibit/.style={-{Bar[width=8pt]}, very thick, draw=bpred,
    shorten >=1pt, shorten <=1pt},
  %% binding / association 结合: plain line, no head (use both ways for
  %% dissociable complexes, or add label "bind")
  bind/.style={draw=black!75, thick, shorten >=1pt, shorten <=1pt},
  %% phosphorylation 磷酸化: arrow + automatic P badge at midpoint
  phosphorylate/.style={-{Triangle[length=6pt]}, thick, draw=bpred!80!black,
    postaction={decorate, decoration={markings,
      mark=at position 0.55 with {\node[pbadge]{P};}}},
    shorten >=1pt, shorten <=1pt},
  %% ubiquitination 泛素化: arrow + Ub badge at midpoint
  ubiquitinate/.style={-{Triangle[length=6pt]}, thick, draw=bppurple,
    postaction={decorate, decoration={markings,
      mark=at position 0.55 with {\node[ubbadge]{Ub};}}},
    shorten >=1pt, shorten <=1pt},
  %% transport / translocation 转运/入核: dashed arrow
  transport/.style={-{Stealth[length=6pt]}, thick, dashed, draw=black!70,
    shorten >=1pt, shorten <=1pt},
  %% recruitment 募集: dotted open arrow
  recruit/.style={-{Stealth[open,length=6pt]}, thick, dotted, draw=bpblue!80,
    shorten >=1pt, shorten <=1pt},
  %% uncertain relationship (无材料依据时): gray dashed + question mark node
  uncertain/.style={-{Stealth[length=6pt]}, thick, dashed, draw=bpgray,
    postaction={decorate, decoration={markings,
      mark=at position 0.5 with {\node[font=\scriptsize\bfseries,
        text=bpgray]{?};}}},
    shorten >=1pt, shorten <=1pt},
}
%% phosphorylate/ubiquitinate/uncertain need the markings library:
\usetikzlibrary{decorations.markings}

%% ---------- legend box ----------
%% A legend is a plain node with the legendframe style containing a tabular
%% of \legendrow{<edge style>}{<text>} rows (see example-mechanism.tex):
%%   \node[legendframe, anchor=north west] at (x,y) {%
%%     \begin{tabular}{@{}ll@{}}
%%       \legendrow{activate}{activation}
%%       \legendrow{inhibit}{inhibition}
%%     \end{tabular}};
%% \legendsymbol{<style>}{<width>} draws one sample stroke for a legend row.
\newcommand{\legendsymbol}[2]{% #1 edge/node style, #2 width
  \tikz[baseline=-0.6ex]{\draw[#1] (0,0) -- (#2,0);}}
\newcommand{\legendrow}[2]{\legendsymbol{#1}{0.7} & #2 \\}
\tikzset{
  legendframe/.style={rectangle, draw=black!60, fill=white, rounded corners,
    inner sep=6pt, align=left, font=\scriptsize},
}

%% ---------- pgfplots chart preamble block ----------
%% Consistent look for data charts drawn from actual CSV/TSV files:
%%   \begin{tikzpicture}\begin{axis}[bpchart, xlabel=..., ylabel=...]
%%     \addplot table[x=group,y=log2FC,col sep=comma] {data.csv};
%%   \end{axis}\end{tikzpicture}
\pgfplotsset{
  bpchart/.style={
    width=9cm, height=6.5cm,
    axis lines=left, enlarge x limits=0.05, enlarge y limits=0.05,
    tick label style={font=\scriptsize}, label style={font=\small},
    legend style={font=\scriptsize, draw=black!50, fill=white,
      rounded corners=1pt},
    grid=major, grid style={black!12},
    every axis plot/.append style={thick},
  },
  bpbar/.style={ybar, bar width=8pt, fill=bpblue!55, draw=bpblue},
  bpscatter/.style={only marks, mark=*, mark size=1.6pt, draw=bpblue,
    fill=bpblue!60},
}
        % this file (same directory or TEXINPUTS)
%%   \begin{document} ... \end{document}
%%
%% Dependencies (deliberately minimal): tikz, pgfplots, xcolor only.
%% TikZ libraries used: arrows.meta, positioning, calc, backgrounds,
%% fit, decorations.pathmorphing -- all ship with every full TeX Live.
%% Compiles with pdflatex (also xelatex). No exotic packages.
%% =====================================================================
\RequirePackage{xcolor}
\RequirePackage{tikz}
\usetikzlibrary{arrows.meta,positioning,calc,backgrounds,fit,decorations.pathmorphing}
\RequirePackage{pgfplots}
\pgfplotsset{compat=1.17} % safe baseline; 1.16 also fine if older TeX Live

%% ---------- palette (colorblind-aware, Okabe-Ito inspired) ----------
\definecolor{bpblue}{HTML}{0072B2}    % protein / receptor fill
\definecolor{bpblueL}{HTML}{BFDFF2}   % protein light fill
\definecolor{bporange}{HTML}{E69F00}  % gene / nucleic fill
\definecolor{bpgreen}{HTML}{009E73}   % small molecule
\definecolor{bppurple}{HTML}{CC79A7}  % complex
\definecolor{bpred}{HTML}{D55E00}     % inhibition / danger
\definecolor{bpgray}{HTML}{999999}    % uncertain
\definecolor{membrane}{HTML}{F2E3C9}  % lipid bilayer band
\definecolor{bpbg}{HTML}{F7F9FB}      % compartment background tint

%% ---------- compartment bands ----------
%% Horizontal bands, top->bottom: extracellular, membrane, cytosol, nucleus.
%% Draw them on the background layer FIRST, then place nodes on top.
%% Usage: \fill[extracellularband] (-8,2.5) rectangle (8,5);
\tikzset{
  extracellularband/.style={fill=white},
  membraneband/.style={fill=membrane},
  cytosolband/.style={fill=bpbg},
  nucleusband/.style={fill=bpblueL!35},
  bandlabel/.style={font=\small\itshape, text=black!55, anchor=north west},
}

%% ---------- membrane bilayer pic ----------
%% A double line of lipid head circles + tails spanning a horizontal band.
%% Usage: \pic at (0,0) {bilayer={8}};  % half-width 8cm
\tikzset{
  pics/bilayer/.style={
    code={
      \pgfmathsetmacro{\hw}{#1}
      \begin{scope}
        \foreach \i in {0,...,40} {
          \pgfmathsetmacro{\px}{-\hw + \i*(2*\hw/40)}
          \draw[fill=bpred!70, draw=bpred!70] (\px,0.28) circle (0.09);
          \draw[fill=bpred!70, draw=bpred!70] (\px,-0.28) circle (0.09);
          \draw[bporange!80, line width=0.7pt] (\px,0.19) -- (\px,0.03);
          \draw[bporange!80, line width=0.7pt] (\px,-0.19) -- (\px,-0.03);
        }
      \end{scope}
    }
  },
}

%% ---------- node styles ----------
\tikzset{
  %% protein: rounded rectangle, upright (正体) label
  protein/.style={rectangle, rounded corners=3pt, draw=bpblue,
    fill=bpblueL, thick, minimum width=1.7cm, minimum height=0.8cm,
    align=center, font=\small},
  %% gene: italic (斜体) label, orange accent. Use \gene{IL6} macro below
  %% to guarantee italic math-style name.
  gene/.style={rectangle, draw=bporange, fill=bporange!18, thick,
    minimum width=1.4cm, minimum height=0.7cm, align=center,
    font=\small\itshape},
  %% complex: two stacked rounded boxes look (purple)
  complex/.style={rectangle, rounded corners=6pt, draw=bppurple,
    fill=bppurple!20, thick, double, double distance=1.2pt,
    minimum width=2.0cm, minimum height=0.9cm, align=center, font=\small},
  %% small molecule / metabolite: ellipse, green
  smallmolecule/.style={ellipse, draw=bpgreen, fill=bpgreen!15, thick,
    minimum width=1.3cm, minimum height=0.7cm, align=center, font=\small},
  %% transmembrane receptor: tall rounded rect meant to SPAN the membrane
  %% band (place its center on the membrane midline).
  receptor/.style={rectangle, rounded corners=8pt, draw=bpblue,
    fill=bpblueL, very thick, minimum width=1.1cm, minimum height=2.0cm,
    align=center, font=\small},
  %% nucleus body (large double circle) for a nuclear compartment
  nucleusbody/.style={circle, draw=bpblue!60, fill=nucleusband, thick,
    double, double distance=2pt, minimum size=5.2cm},
  %% small badges: phosphate / ubiquitin
  pbadge/.style={circle, draw=bpred, fill=white, inner sep=0.6pt,
    minimum size=0.42cm, font=\scriptsize\bfseries},
  ubbadge/.style={circle, draw=bppurple, fill=white, inner sep=0.6pt,
    minimum size=0.55cm, font=\scriptsize\bfseries},
  %% generic annotation text
  ann/.style={font=\scriptsize, text=black!70, align=center},
}

%% Italic gene name helper: \gene{STAT3} -> italic STAT3
\newcommand{\gene}[1]{\textit{#1}}
%% Phosphorylated protein: \phos{STAT3} -> p-STAT3 (upright p- prefix)
\newcommand{\phos}[1]{p-#1}

%% ---------- edge styles ----------
%% arrows.meta tips; every edge carries its mechanism semantics in style.
\tikzset{
  %% activation 促进/激活: open triangle arrow
  activate/.style={-{Triangle[open,length=7pt]}, thick, draw=bpblue,
    shorten >=1pt, shorten <=1pt},
  %% inhibition 抑制: perpendicular T-bar
  inhibit/.style={-{Bar[width=8pt]}, very thick, draw=bpred,
    shorten >=1pt, shorten <=1pt},
  %% binding / association 结合: plain line, no head (use both ways for
  %% dissociable complexes, or add label "bind")
  bind/.style={draw=black!75, thick, shorten >=1pt, shorten <=1pt},
  %% phosphorylation 磷酸化: arrow + automatic P badge at midpoint
  phosphorylate/.style={-{Triangle[length=6pt]}, thick, draw=bpred!80!black,
    postaction={decorate, decoration={markings,
      mark=at position 0.55 with {\node[pbadge]{P};}}},
    shorten >=1pt, shorten <=1pt},
  %% ubiquitination 泛素化: arrow + Ub badge at midpoint
  ubiquitinate/.style={-{Triangle[length=6pt]}, thick, draw=bppurple,
    postaction={decorate, decoration={markings,
      mark=at position 0.55 with {\node[ubbadge]{Ub};}}},
    shorten >=1pt, shorten <=1pt},
  %% transport / translocation 转运/入核: dashed arrow
  transport/.style={-{Stealth[length=6pt]}, thick, dashed, draw=black!70,
    shorten >=1pt, shorten <=1pt},
  %% recruitment 募集: dotted open arrow
  recruit/.style={-{Stealth[open,length=6pt]}, thick, dotted, draw=bpblue!80,
    shorten >=1pt, shorten <=1pt},
  %% uncertain relationship (无材料依据时): gray dashed + question mark node
  uncertain/.style={-{Stealth[length=6pt]}, thick, dashed, draw=bpgray,
    postaction={decorate, decoration={markings,
      mark=at position 0.5 with {\node[font=\scriptsize\bfseries,
        text=bpgray]{?};}}},
    shorten >=1pt, shorten <=1pt},
}
%% phosphorylate/ubiquitinate/uncertain need the markings library:
\usetikzlibrary{decorations.markings}

%% ---------- legend box ----------
%% A legend is a plain node with the legendframe style containing a tabular
%% of \legendrow{<edge style>}{<text>} rows (see example-mechanism.tex):
%%   \node[legendframe, anchor=north west] at (x,y) {%
%%     \begin{tabular}{@{}ll@{}}
%%       \legendrow{activate}{activation}
%%       \legendrow{inhibit}{inhibition}
%%     \end{tabular}};
%% \legendsymbol{<style>}{<width>} draws one sample stroke for a legend row.
\newcommand{\legendsymbol}[2]{% #1 edge/node style, #2 width
  \tikz[baseline=-0.6ex]{\draw[#1] (0,0) -- (#2,0);}}
\newcommand{\legendrow}[2]{\legendsymbol{#1}{0.7} & #2 \\}
\tikzset{
  legendframe/.style={rectangle, draw=black!60, fill=white, rounded corners,
    inner sep=6pt, align=left, font=\scriptsize},
}

%% ---------- pgfplots chart preamble block ----------
%% Consistent look for data charts drawn from actual CSV/TSV files:
%%   \begin{tikzpicture}\begin{axis}[bpchart, xlabel=..., ylabel=...]
%%     \addplot table[x=group,y=log2FC,col sep=comma] {data.csv};
%%   \end{axis}\end{tikzpicture}
\pgfplotsset{
  bpchart/.style={
    width=9cm, height=6.5cm,
    axis lines=left, enlarge x limits=0.05, enlarge y limits=0.05,
    tick label style={font=\scriptsize}, label style={font=\small},
    legend style={font=\scriptsize, draw=black!50, fill=white,
      rounded corners=1pt},
    grid=major, grid style={black!12},
    every axis plot/.append style={thick},
  },
  bpbar/.style={ybar, bar width=8pt, fill=bpblue!55, draw=bpblue},
  bpscatter/.style={only marks, mark=*, mark size=1.6pt, draw=bpblue,
    fill=bpblue!60},
}
        % this file (same directory or TEXINPUTS)
%%   \begin{document} ... \end{document}
%%
%% Dependencies (deliberately minimal): tikz, pgfplots, xcolor only.
%% TikZ libraries used: arrows.meta, positioning, calc, backgrounds,
%% fit, decorations.pathmorphing -- all ship with every full TeX Live.
%% Compiles with pdflatex (also xelatex). No exotic packages.
%% =====================================================================
\RequirePackage{xcolor}
\RequirePackage{tikz}
\usetikzlibrary{arrows.meta,positioning,calc,backgrounds,fit,decorations.pathmorphing}
\RequirePackage{pgfplots}
\pgfplotsset{compat=1.17} % safe baseline; 1.16 also fine if older TeX Live

%% ---------- palette (colorblind-aware, Okabe-Ito inspired) ----------
\definecolor{bpblue}{HTML}{0072B2}    % protein / receptor fill
\definecolor{bpblueL}{HTML}{BFDFF2}   % protein light fill
\definecolor{bporange}{HTML}{E69F00}  % gene / nucleic fill
\definecolor{bpgreen}{HTML}{009E73}   % small molecule
\definecolor{bppurple}{HTML}{CC79A7}  % complex
\definecolor{bpred}{HTML}{D55E00}     % inhibition / danger
\definecolor{bpgray}{HTML}{999999}    % uncertain
\definecolor{membrane}{HTML}{F2E3C9}  % lipid bilayer band
\definecolor{bpbg}{HTML}{F7F9FB}      % compartment background tint

%% ---------- compartment bands ----------
%% Horizontal bands, top->bottom: extracellular, membrane, cytosol, nucleus.
%% Draw them on the background layer FIRST, then place nodes on top.
%% Usage: \fill[extracellularband] (-8,2.5) rectangle (8,5);
\tikzset{
  extracellularband/.style={fill=white},
  membraneband/.style={fill=membrane},
  cytosolband/.style={fill=bpbg},
  nucleusband/.style={fill=bpblueL!35},
  bandlabel/.style={font=\small\itshape, text=black!55, anchor=north west},
}

%% ---------- membrane bilayer pic ----------
%% A double line of lipid head circles + tails spanning a horizontal band.
%% Usage: \pic at (0,0) {bilayer={8}};  % half-width 8cm
\tikzset{
  pics/bilayer/.style={
    code={
      \pgfmathsetmacro{\hw}{#1}
      \begin{scope}
        \foreach \i in {0,...,40} {
          \pgfmathsetmacro{\px}{-\hw + \i*(2*\hw/40)}
          \draw[fill=bpred!70, draw=bpred!70] (\px,0.28) circle (0.09);
          \draw[fill=bpred!70, draw=bpred!70] (\px,-0.28) circle (0.09);
          \draw[bporange!80, line width=0.7pt] (\px,0.19) -- (\px,0.03);
          \draw[bporange!80, line width=0.7pt] (\px,-0.19) -- (\px,-0.03);
        }
      \end{scope}
    }
  },
}

%% ---------- node styles ----------
\tikzset{
  %% protein: rounded rectangle, upright (正体) label
  protein/.style={rectangle, rounded corners=3pt, draw=bpblue,
    fill=bpblueL, thick, minimum width=1.7cm, minimum height=0.8cm,
    align=center, font=\small},
  %% gene: italic (斜体) label, orange accent. Use \gene{IL6} macro below
  %% to guarantee italic math-style name.
  gene/.style={rectangle, draw=bporange, fill=bporange!18, thick,
    minimum width=1.4cm, minimum height=0.7cm, align=center,
    font=\small\itshape},
  %% complex: two stacked rounded boxes look (purple)
  complex/.style={rectangle, rounded corners=6pt, draw=bppurple,
    fill=bppurple!20, thick, double, double distance=1.2pt,
    minimum width=2.0cm, minimum height=0.9cm, align=center, font=\small},
  %% small molecule / metabolite: ellipse, green
  smallmolecule/.style={ellipse, draw=bpgreen, fill=bpgreen!15, thick,
    minimum width=1.3cm, minimum height=0.7cm, align=center, font=\small},
  %% transmembrane receptor: tall rounded rect meant to SPAN the membrane
  %% band (place its center on the membrane midline).
  receptor/.style={rectangle, rounded corners=8pt, draw=bpblue,
    fill=bpblueL, very thick, minimum width=1.1cm, minimum height=2.0cm,
    align=center, font=\small},
  %% nucleus body (large double circle) for a nuclear compartment
  nucleusbody/.style={circle, draw=bpblue!60, fill=nucleusband, thick,
    double, double distance=2pt, minimum size=5.2cm},
  %% small badges: phosphate / ubiquitin
  pbadge/.style={circle, draw=bpred, fill=white, inner sep=0.6pt,
    minimum size=0.42cm, font=\scriptsize\bfseries},
  ubbadge/.style={circle, draw=bppurple, fill=white, inner sep=0.6pt,
    minimum size=0.55cm, font=\scriptsize\bfseries},
  %% generic annotation text
  ann/.style={font=\scriptsize, text=black!70, align=center},
}

%% Italic gene name helper: \gene{STAT3} -> italic STAT3
\newcommand{\gene}[1]{\textit{#1}}
%% Phosphorylated protein: \phos{STAT3} -> p-STAT3 (upright p- prefix)
\newcommand{\phos}[1]{p-#1}

%% ---------- edge styles ----------
%% arrows.meta tips; every edge carries its mechanism semantics in style.
\tikzset{
  %% activation 促进/激活: open triangle arrow
  activate/.style={-{Triangle[open,length=7pt]}, thick, draw=bpblue,
    shorten >=1pt, shorten <=1pt},
  %% inhibition 抑制: perpendicular T-bar
  inhibit/.style={-{Bar[width=8pt]}, very thick, draw=bpred,
    shorten >=1pt, shorten <=1pt},
  %% binding / association 结合: plain line, no head (use both ways for
  %% dissociable complexes, or add label "bind")
  bind/.style={draw=black!75, thick, shorten >=1pt, shorten <=1pt},
  %% phosphorylation 磷酸化: arrow + automatic P badge at midpoint
  phosphorylate/.style={-{Triangle[length=6pt]}, thick, draw=bpred!80!black,
    postaction={decorate, decoration={markings,
      mark=at position 0.55 with {\node[pbadge]{P};}}},
    shorten >=1pt, shorten <=1pt},
  %% ubiquitination 泛素化: arrow + Ub badge at midpoint
  ubiquitinate/.style={-{Triangle[length=6pt]}, thick, draw=bppurple,
    postaction={decorate, decoration={markings,
      mark=at position 0.55 with {\node[ubbadge]{Ub};}}},
    shorten >=1pt, shorten <=1pt},
  %% transport / translocation 转运/入核: dashed arrow
  transport/.style={-{Stealth[length=6pt]}, thick, dashed, draw=black!70,
    shorten >=1pt, shorten <=1pt},
  %% recruitment 募集: dotted open arrow
  recruit/.style={-{Stealth[open,length=6pt]}, thick, dotted, draw=bpblue!80,
    shorten >=1pt, shorten <=1pt},
  %% uncertain relationship (无材料依据时): gray dashed + question mark node
  uncertain/.style={-{Stealth[length=6pt]}, thick, dashed, draw=bpgray,
    postaction={decorate, decoration={markings,
      mark=at position 0.5 with {\node[font=\scriptsize\bfseries,
        text=bpgray]{?};}}},
    shorten >=1pt, shorten <=1pt},
}
%% phosphorylate/ubiquitinate/uncertain need the markings library:
\usetikzlibrary{decorations.markings}

%% ---------- legend box ----------
%% A legend is a plain node with the legendframe style containing a tabular
%% of \legendrow{<edge style>}{<text>} rows (see example-mechanism.tex):
%%   \node[legendframe, anchor=north west] at (x,y) {%
%%     \begin{tabular}{@{}ll@{}}
%%       \legendrow{activate}{activation}
%%       \legendrow{inhibit}{inhibition}
%%     \end{tabular}};
%% \legendsymbol{<style>}{<width>} draws one sample stroke for a legend row.
\newcommand{\legendsymbol}[2]{% #1 edge/node style, #2 width
  \tikz[baseline=-0.6ex]{\draw[#1] (0,0) -- (#2,0);}}
\newcommand{\legendrow}[2]{\legendsymbol{#1}{0.7} & #2 \\}
\tikzset{
  legendframe/.style={rectangle, draw=black!60, fill=white, rounded corners,
    inner sep=6pt, align=left, font=\scriptsize},
}

%% ---------- pgfplots chart preamble block ----------
%% Consistent look for data charts drawn from actual CSV/TSV files:
%%   \begin{tikzpicture}\begin{axis}[bpchart, xlabel=..., ylabel=...]
%%     \addplot table[x=group,y=log2FC,col sep=comma] {data.csv};
%%   \end{axis}\end{tikzpicture}
\pgfplotsset{
  bpchart/.style={
    width=9cm, height=6.5cm,
    axis lines=left, enlarge x limits=0.05, enlarge y limits=0.05,
    tick label style={font=\scriptsize}, label style={font=\small},
    legend style={font=\scriptsize, draw=black!50, fill=white,
      rounded corners=1pt},
    grid=major, grid style={black!12},
    every axis plot/.append style={thick},
  },
  bpbar/.style={ybar, bar width=8pt, fill=bpblue!55, draw=bpblue},
  bpscatter/.style={only marks, mark=*, mark size=1.6pt, draw=bpblue,
    fill=bpblue!60},
}


\begin{document}
\begin{tikzpicture}[node distance=1.1cm and 1.6cm]

  %% ---- compartment bands (background first) ----
  \begin{scope}[on background layer]
    \fill[extracellularband] (-6.5, 1.4) rectangle (6.5, 3.4);
    \fill[membraneband]      (-6.5, 0.6) rectangle (6.5, 1.4);
    \fill[cytosolband]       (-6.5,-3.6) rectangle (6.5, 0.6);
    \fill[nucleusband]       ( 1.0,-3.6) rectangle (6.5,-1.2);
  \end{scope}
  \node[bandlabel] at (-6.4, 3.35) {extracellular};
  \node[bandlabel] at (-6.4, 1.38) {membrane};
  \node[bandlabel] at (-6.4, 0.55) {cytosol};
  \node[bandlabel] at ( 1.1,-1.25) {nucleus};
  \pic at (0,1.0) {bilayer={6.5}};

  %% ---- entities ----
  \node[smallmolecule] (lig) at (-3.5, 2.4) {LIGAND-X};
  \node[receptor]      (rcp) at (-3.5, 1.0) {RcptA};
  \node[protein]       (kin) at (-1.2,-0.8) {KinB};
  \node[protein]       (rep) at (-4.6,-2.2) {RepE};
  \node[protein]       (tf)  at ( 3.6,-1.9) {TFc};
  \node[gene]          (tgt) at ( 3.6,-3.0) {\gene{TargetD}};

  %% ---- edges (mechanism semantics) ----
  \draw[activate]      (lig) -- (rcp);
  \draw[activate]      (rcp.south) |- (kin.west);
  \draw[phosphorylate] (kin) -- (tf);
  \draw[transport]     (kin.south) -- ++(0,-0.9) node[ann, right] {shuttle};
  \draw[inhibit]       (rep) -- (kin);
  \draw[activate]      (tf) -- (tgt);

  %% ---- legend ----
  \node[legendframe, anchor=north west] at (-6.4, -3.8) {%
    \begin{tabular}{@{}ll@{}}
      \multicolumn{2}{@{}l}{\textbf{Legend (synthetic)}}\\[2pt]
      \legendrow{activate}{activation}%
      \legendrow{inhibit}{inhibition}%
      \legendrow{phosphorylate}{phosphorylation}%
    \end{tabular}};

\end{tikzpicture}
\end{document}
